\documentclass[11pt]{article}

\usepackage[a4paper,margin=1in]{geometry}
\usepackage[T1]{fontenc}
\usepackage[utf8]{inputenc}
\usepackage{lmodern}
\usepackage{microtype}
\usepackage{setspace}
\usepackage{parskip}
\usepackage{hyperref}
\usepackage{enumitem}

\onehalfspacing

\title{Constitutional AI under Finite Capacity\\
Instrument Index}
\author{Panagiotis Kalomoirakis\\
Independent Researcher, Synkyria Project}
\date{April 2026}

\begin{document}

\maketitle


\begin{center}
\small\textsc{Canonical Release v1.0}

\small \textcopyright\ Panagiotis Kalomoirakis 2026. All rights reserved.

\end{center}

\noindent This index is included in order to make the structure of the Constitutional AI under Finite Capacity package explicit. The package is organised as a constitutional core together with companion instruments that clarify force, projection, evidence, interpretive limits, and public-governance translation.

\section*{Constituent instruments}

\begin{enumerate}[leftmargin=2em]
    \item \textbf{The Finite-Capacity AI Constitution}  
    The constitutional core. It states the minimal structural articles required for accountable AI operation under finite capacity.

    \item \textbf{Constitutional Force and Applicability in the Finite-Capacity AI Constitution}  
    An interpretive companion clarifying how the constitutional articles should be read canonically, including universal minimum force, governance-strengthening clauses, and profile-dependent historical qualification.

    \item \textbf{Minimal Application Profile for Constitutional AI in AI Assistant Systems}  
    A projection instrument showing how the constitutional grammar is applied to AI assistant systems.

    \item \textbf{Minimal Reviewable Evidence Specification for Constitutional AI}  
    A review instrument specifying the minimum evidence surface required for constitutional claims to remain inspectable.

    \item \textbf{Questions, Limits, and Misreadings for Constitutional AI under Finite Capacity}  
    A clarification instrument addressing predictable questions, limits, and interpretive errors.

    \item \textbf{Constitutional AI under Finite Capacity: A Policy Translation Brief}  
    A public-governance translation instrument intended for policy audiences, public authorities, compliance teams, standards-facing groups, and decision-makers. It translates the constitutional package into outward-facing governance language without altering the constitutional core.
\end{enumerate}

\section*{Recommended reading order}

\begin{enumerate}[leftmargin=2em]
    \item The Finite-Capacity AI Constitution
    \item Constitutional Force and Applicability in the Finite-Capacity AI Constitution
    \item Minimal Application Profile for Constitutional AI in AI Assistant Systems
    \item Minimal Reviewable Evidence Specification for Constitutional AI
    \item Questions, Limits, and Misreadings for Constitutional AI under Finite Capacity
    \item Constitutional AI under Finite Capacity: A Policy Translation Brief
\end{enumerate}

\noindent For internal reading, the order above should be preserved. For outward-facing release, however, the \emph{Policy Translation Brief} may be presented first as an entry document for policy, regulatory, public-authority, standards, and compliance audiences, while the constitutional core remains the primary object of the package.

\section*{Interpretive note}

The constitutional core should be read first and treated as the primary object. The remaining documents do not replace it. They clarify how it should be read, how it is projected into a concrete system class, what review surface is required for its claims, how likely misunderstandings should be handled without overburdening the constitutional text itself, and how the package may be translated into outward-facing public-governance language.

\section*{Scope note}

This package is not a single implementation manual, a complete ethics theory, or a legal code. It is a structured constitutional and governance-oriented series concerning accountable AI under finite capacity.

\end{document}